\documentclass[12pt]{article}

% --- Packages ---
\usepackage[margin=1in]{geometry}% Sets standard 1-inch margins
\usepackage{amsmath,amssymb}% Standard math packages
\usepackage{graphicx}% Required for inserting images
\usepackage{float}% Allows precise image placement using [H]
\usepackage{booktabs}% For professional-looking tables
\usepackage{hyperref}% Makes references and citations clickable
\usepackage{setspace}% Allows line spacing adjustments
\usepackage{caption}% Customizing captions for figures/tables

% --- Setup hyperref ---
\hypersetup{
    colorlinks=true,
    linkcolor=blue,
    filecolor=magenta,
    urlcolor=cyan,
    citecolor=blue,
}

% --- Title Information ---
\title{\textbf{Final Project Report:\\How Music Taste Changes Over Time}}
\author{
    Lauren Yee\\
    Kian Sharifzadeh\\
    Krish Doshi\\
    \vspace{0.3cm}
    \textit{Stats 423: Applied Linear Regression $\&$ ANOVA}
}
\date{\today}

\begin{document}

\maketitle
\begin{abstract}
\noindent Music is a key part of daily life, yet people's opinions about music have evolved over time. 
\noindent With factors, such as race, having a different impact on society during each decade, this study focuses on
\noindent how various factors affect song ratings during different decades. Data from Billboard Hot 100 Number One songs were
\noindent taken from the 1960s - 2010s and split up to form rating prediction linear models. We compared main effects models
\noindent with interaction effects models and decided which one would be more viable to optimize. The chosen model for the decade
\noindent was optimized based on its AIC score to determine which predictors made a significant impact on the model. Comparing the predictors across
\noindent models revealed that variables such as energy and race were consistently important across most decades. 
\noindent However, no singular factor was consistent between every decade. Ultimately, the findings demonstrated that 
\noindent while the factors that predict song ratings change over time, they have the potential to be similar in different time periods.
\end{abstract}

\onehalfspacing % 1.5 spacing makes it easier for grading
 
% -------------------------------------------------------------------
\section{Introduction}
% -------------------------------------------------------------------
Introduce your research question and the motivation behind your study. Why is this topic interesting or important? State the primary objectives of your regression analysis. 

We have 3 main research questions. They include: 

    \begin{enumerate}
    \item  How do demographic characteristics of the artist affect the overall rating of a song over time?
    
    \item How do song characteristics (provided by Spotify) affect the overall rating of a song over time?
    
   \item  Which demographics and song characteristics significantly impact the overall rating of a Billboard \#1 song over the 6 decades?
\end{enumerate}

We were interested in how music has changed over time and we thought separating our data into different time periods helps us 
understand what characteristics truly impact the overall rating. We thought each decade had its own music taste, so we decided 
to analyze the data based on each decade the hits were created.
It is interesting to see how far music has gone over the years, as well as getting a better grasp of music taste in each decade.

% -------------------------------------------------------------------
\section{Data Description and Exploratory Data Analysis (EDA)}
% -------------------------------------------------------------------
Describe the dataset you are using. Where did it come from? What are the observational units? 

\subsection{Variables}
Briefly explain the response variable ($Y$) and the main predictor variables ($X_i$). 

% Example of an inline equation and bullet points
Let $Y$ denote the overall rating of the song. The primary predictors are:
\begin{itemize}
    \item $X_1$: Age: Average age of the artist(s)
    \item $X_2$: Gender: Gender of the artist (s)
    \item $X_3$: Race: Race of the artist (s)
    \item $X_4$: Genre: Type of genre of the song
    \item $X_5$: Energy: Energy measure from Spotify (0-100) 
    \item $X_6$: Danceability: Danceability measure from Spotify (0-100) 
    \item $X_7$: Happiness: Happiness measure from Spotify (0-100)
    \item $X_8$: Loudness_db: Loudness measured in decibels from Spotify
    \item $X_9$: BPM: Beats per minute from Spotify
    \item $X_10$: Length_sec: Length of the song in seconds
\end{itemize}

\subsection{Exploratory Analysis}
Present initial findings from your EDA. You can reference figures like scatterplots or correlation matrices here (e.g., "As seen in Figure \ref{fig:scatter}...").

% Example of inserting a figure
\begin{figure}[htbp] 
    \centering
    \includegraphics[angle=270, width=0.7\textwidth]{IMG_2679.jpeg}
    \caption{A descriptive caption explaining the plot.}
    \label{fig:scatter} % Use this label to reference the figure in your text
\end{figure}

% -------------------------------------------------------------------
\section{Methodology and Model Building}
% -------------------------------------------------------------------
Describe the models you fit to the data. 

\subsection{Proposed Model}
% Example of a standalone equation
We fit a multiple linear regression model of the form:
\begin{equation}
    Y_i = \beta_0 + \beta_1 X_{i1} + \beta_2 X_{i2} + \dots + \beta_p X_{ip} + \epsilon_i
    \label{eq:main_model}
\end{equation}
where $\epsilon_i \overset{iid}{\sim} \mathcal{N}(0, \sigma^2)$.

Explain how you selected your variables (e.g., stepwise selection, theoretical justification, handling of multicollinearity). If you used transformations (like taking the log of the response), explain why.

% -------------------------------------------------------------------
\section{Results and Diagnostics}
% -------------------------------------------------------------------
Present the results of your final model. 

\subsection{Model Summary}
% Example of a professional table
Table \ref{tab:results} summarizes the estimated coefficients for our final model.

\begin{table}[H]
    \centering
    \caption{Estimated Regression Coefficients and Significance}
    \vspace{0.2cm} % Adds a little space between caption and table
    \label{tab:results}
    \begin{tabular}{lrlrr}
        \toprule
        \textbf{Predictor} & \textbf{Estimate} & \textbf{Std. Error} & \textbf{t-value} & \textbf{Pr($>|t|$)} \\
        \midrule
        (Intercept) & $12.45$ & $2.10$ & $5.93$ & $<0.001$ \\
        $X_1$       & $3.12$  & $0.45$ & $6.93$ & $<0.001$ \\
        $X_2$       & $-1.05$ & $0.30$ & $-3.50$ & $0.002$ \\
        \bottomrule
    \end{tabular}
\end{table}

\subsection{Model Diagnostics}
Assess the assumptions of linear regression (linearity, independence, homoscedasticity, normality of residuals). You should include residual plots or QQ-plots here to justify that your model assumptions hold (or discuss how they are violated).

% -------------------------------------------------------------------
\section{Discussion and Conclusion}
% -------------------------------------------------------------------
Summarize the main takeaways from your model in the context of the original problem. What do the coefficients mean in the real world? 

Acknowledge the limitations of your study (e.g., potential omitted variable bias, non-representative sample) and suggest areas for future work.

% -------------------------------------------------------------------
\section*{References}
% -------------------------------------------------------------------
% Using a simple thebibliography environment for ease of use. 
% For larger projects, students might prefer BibTeX.
\begin{thebibliography}{9}
\bibitem{textbook} 
Author Name. 
\textit{Title of the Textbook}. 
Publisher, Year.

\bibitem{dataset} 
Dataset Creator. 
``Title of the Dataset." 
\textit{Source/URL}, Year.
\end{thebibliography}

% -------------------------------------------------------------------
\newpage
\appendix
\section{Appendix: R/Python Code}
% -------------------------------------------------------------------
Paste your analysis code here. If you prefer, you can use the \texttt{verbatim} environment to format it as raw text:

\begin{verbatim}
# Example R Code
model <- lm(Y ~ X1 + X2, data = mydata)
summary(model)
plot(model)
\end{verbatim}

\end{document}